\subsection{Product of a space by a simply connected space}

PROPOSITION 8. --- Let B be a topological space. Let T be a simply connected and locally connected space. Let E be a covering of B × T, with projection p, and let t be a point of T. Denote by \(E_{t}\) the space \(p^{-1}(B \times \{t\})\) equipped with the map \(p_{t} = pr_{1} \circ p \mid E_{t}: E_{t} \to B\) ; this is a covering of B. There then exists a unique (B × T)-isomorphism from the covering ( \(E_{t} \times T, p_{t} \times Id_{T}\) ) onto the covering E which maps (x, t) to x for all \(x \in E_{t}\) .

One may assume that B is nonempty. Let \(x\) be a point of \(E_t\) . By Proposition 3 of I, p. 125 applied to the covering E and the continuous map T → B × T, \(u \mapsto (p_t(x), u)\) , there exists a unique continuous map \(f_x \colon \mathrm{T} \to \mathrm{E}\) such that \(f_x(t) = x\) and \(p(f_x(u)) = (p_t(x), u)\) for all \(u \in T\) . Let \(h \colon E_t \times T \to E\) be the map defined by \(h(x, u) = f_x(u)\) . We have \(h(x, t) = x\) and \(p \circ h = p_t \times \operatorname{Id}_T\) . The map h is a lift to E of the map \(p_t \times \operatorname{Id}_T\) . The restriction of h to \(E_t \times \{t\}\) is continuous, as is the restriction of h to \(\{x\} \times T\) for every point x of \(E_t\) . Since the space T is locally connected, the map h is continuous (I, p. 37, cor. 1 of th. 1).

Let \(b\) be a point of B. By construction, the map \(h\) induces a bijection from the fiber \(p_t^{-1}(b) \times \{t\}\) of \(E_t \times T\) onto the fiber \(p^{-1}(b,t)\) of E at \((b,t)\) . Since the space T is connected and locally connected, the map \(h\) is bijective (I, p. 84, cor. of prop. 7). It is therefore a \(B \times T\) -isomorphism (I, p. 30, cor. 2 of prop. 6).

Let \(h'\) be a \((\mathrm{B} \times \mathrm{T})\) -isomorphism from the covering \((\mathrm{E}_{t} \times \mathrm{T}, p_{t} \times \mathrm{Id}_{\mathrm{T}})\) onto the covering E which maps \((x, t)\) to x for every point \(x \in E_{t}\) . For all \(x \in E_{t}\) , the maps \(u \mapsto h(x, u)\) and \(u \mapsto h'(x, u)\) are equal (I, p. 34, cor. 1 of prop. 11). We therefore have \(h = h'\) .

COROLLARY 1. --- Under the hypotheses of Prop. 8, if t and t' are two points of T, the coverings \(E_{t}\) and \(E_{t'}\) are B-isomorphic.

COROLLARY 2. --- Under the hypotheses of Prop. 8, let (E,p) and (E',p') be coverings of B × T and let \(k: E_t \to E_t'\) be a B-morphism. There exists a unique (B × T)-morphism \(\widetilde{k}\) : E → E' which extends k. If k is a B-isomorphism, \(\widetilde{k}\) is a (B × T)-isomorphism.

By proposition 8, we may assume that there exist coverings F and F′ of B such that E and E′ are respectively the (B × T)-spaces F × T and F′ × T. We then have \(E_{t} = F \times \{t\}\) , \(E_{t}' = F' \times \{t\}\) and the map k is written \((x, t) \mapsto (k'(x), t)\) , where \(k'\) is a B-morphism from F into F'. The map \(k' \times Id_{T}\) is a B × T-morphism extending k; it is an isomorphism if k is one.

Let \(\widetilde{k}\colon\mathrm{E}\to\mathrm{E}^{\prime}\) be a \((\mathrm{B}\times\mathrm{T})\) -morphism extending \(k\) . Let \(x\) be a point of F. Denote by \(q\) the projection of F and set \(b=q(x)\) . The maps \(u\mapsto\widetilde{k}(x,u)\) and \(u\mapsto(k^{\prime}(x),u)\) are liftings in \(\mathrm{E}^{\prime}\) of the map \(u\mapsto(b,u)\) from T to \(\mathrm{B}\times\mathrm{T}\) . They coincide at \(t\) . Since the space T is connected, they are equal. Thus, \(\widetilde{k}\) is the \((\mathrm{B}\times\mathrm{T})\) -morphism \(k^{\prime}\times\mathrm{Id}_{\mathrm{T}}\colon\mathrm{F}\times\mathrm{T}\to\mathrm{F}^{\prime}\times\mathrm{T}\) .

COROLLARY 3. --- Let B and B' be topological spaces, let T be a simply connected and locally connected topological space. Let \(f: B' \times T \to B\) be a continuous map and let E be a covering of B. Given a point t of T and a continuous lifting \(g_t: B' \times \{t\} \to E\) of \(f | B' \times \{t\}\) to E, there exists a unique continuous lifting g of f to E extending \(g_t\) .

In view of prop. 3 of I, p. 9, it suffices to show that every continuous section of \(f^{*}(\mathrm{E})\) over \(\mathrm{B}' \times \{t\}\) extends uniquely to a continuous section of \(f^{*}(\mathrm{E})\) . Now this follows from cor. 2 applied to the coverings \((\mathrm{B}' \times \mathrm{T}, \mathrm{Id}_{\mathrm{B}' \times \mathrm{T}})\) and \(f^{*}\mathrm{E}\) of \(\mathrm{B}' \times \mathrm{T}\) .

Remark. --- Keep the hypotheses and notation of Proposition 8. Let G be a discrete topological group. Suppose that E is a principal covering with group G. If one equips the coverings \(E_{t}\) of B and \(E_{t} \times T\) of \(B \times T\) with principal covering structures of group G deduced from that of E (I, p. 92, examples 1 and 4), then the coverings E and \(E_{t} \times T\) are isomorphic principal coverings. Indeed, let \(h: E_{t} \times T \to E\) be the unique (B × T)-morphism such that \(h(x, t) = x\) for all \(x \in E_{t}\) . For every element g of G, the map \((x, u) \mapsto h(x \cdot g, u) \cdot g^{-1}\) is a (B × T)-morphism that maps \((x, t)\) to x for all \(x \in E_{t}\) and therefore is equal to h. This proves that h is a (B × T)-morphism of principal coverings.

In particular, under the preceding hypotheses, the principal coverings \(E_{t}\) and \(E_{t'}\) are isomorphic, for \(t' \in T\) . Similarly, one proves that if, in Corollary 2, E and \(E'\) are principal coverings with group G and if k is a B-isomorphism of principal coverings with group G, \(\widetilde{k}\) is a \((\mathrm{B} \times \mathrm{T})\) isomorphism of principal coverings.

PROPOSITION 9. --- Let B be a topological space and let (E,p) be a covering of B. Let T be a simply connected, locally connected, and locally compact topological space. Denote by \(\widetilde{p}\colon\mathcal{C}_{c}(\mathrm{T};\mathrm{E})\to\mathcal{C}_{c}(\mathrm{T};\mathrm{B})\) the map \(g\mapsto p\circ g\) . Let t be a point of T. Denote \(e_{\mathrm{E}}\colon\mathcal{C}_{c}(\mathrm{T};\mathrm{E})\to\mathrm{E}\) the map which to \(g\in\mathcal{C}_{c}(\mathrm{T};\mathrm{E})\) associates \(g(t)\) , and \(e_{\mathrm{B}}\colon\mathcal{C}_{c}(\mathrm{T};\mathrm{B})\to\mathrm{B}\) the map which to \(f\in\mathcal{C}_{c}(\mathrm{T};\mathrm{B})\) associates \(f(t)\) .

The square

\[
\begin{array}{c} \mathcal {C} _ {c} (\mathrm{T}; \mathrm{E}) \xrightarrow {e _ {\mathrm{E}}} \mathrm{E} \\ \Big \downarrow \widetilde {p} \\ \mathcal {C} _ {c} (\mathrm{T}; \mathrm{B}) \xrightarrow {e _ {\mathrm{B}}} \mathrm{B} \end{array}
\]

is Cartesian.

Let us first prove a lemma.

Lemma. --- Let X, Y, Z be topological spaces and let \(g: Y \rightarrow Z\) be a continuous map.

\begin{enumerate}
\def\labelenumi{\alph{enumi})}
\item
  The map \(f\mapsto g\circ f\) from \(\mathcal{C}_c(\mathrm{X};\mathrm{Y})\) to \(\mathcal{C}_c(\mathrm{X};\mathrm{Z})\) is continuous.
\item
  If the topological space Z is Hausdorff, the map \(h \mapsto h \circ g\) from \(\mathcal{C}_{c}(\mathrm{Z};\mathrm{X})\) to \(\mathcal{C}_{c}(\mathrm{Y};\mathrm{X})\) is continuous.
\end{enumerate}

Given a compact subset K of X, an open subset U of Z, and a continuous map f from X to Y, in order to have \((g \circ f)(\mathrm{K}) \subset \mathrm{U}\) , it is necessary and sufficient that one have \(f(\mathrm{K}) \subset g^{-1}(\mathrm{U})\) . The first assertion therefore follows from the definition of the topology of compact convergence (TG, X, p. 26, def. 1).

Similarly, let K be a compact subset of Y and U an open subset of X. Since the space Z is assumed to be Hausdorff, the set \(g(\mathrm{K})\) is compact (TG, I, p. 63, cor. 1). If h is a mapping from Z into X, the condition \((h \circ g)(\mathrm{K}) \subset \mathrm{U}\) is none other than the condition \(h(g(\mathrm{K})) \subset \mathrm{U}\) , whence the second assertion.

Let us now prove Proposition 9. By the lemma, the map \(\widetilde{p}\) is continuous; by Remark 1 of TG, X, p. 27, the maps \(e_{\mathrm{E}}\) and \(e_{\mathrm{B}}\) are continuous. For any map \(g \in \mathcal{C}_c(\mathrm{T};\mathrm{E})\) , we have

\[
(p \circ e _ {\mathrm{E}}) (g) = p (g (t)) = (p \circ g) (t) = \widetilde {p} (g) (t) = \left(e _ {\mathrm{B}} \circ \widetilde {p}\right) (g),
\]

so that the square diagram of the proposition is commutative.

Let \(\varphi\colon\mathcal{C}_{c}(\mathrm{T};\mathrm{E})\to\mathcal{C}_{c}(\mathrm{T};\mathrm{B})\times_{\mathrm{B}}\mathrm{E}\) be the continuous map defined by \(\varphi(g)=(p\circ g,g(t))\) for all \(g\in\mathcal{C}_{c}(\mathrm{T};\mathrm{E})\) . By Proposition 3 of I, p. 125, it is bijective. Indeed, for every pair \((f,x)\in\mathcal{C}_{c}(\mathrm{T};\mathrm{B})\times_{\mathrm{B}}\mathrm{E},\varphi^{-1}(f,x)\) is the unique continuous lift \(g\) of \(f\) to E such that \(g(t)=x\) . Since the space T is locally compact, the map \(\psi\colon(\mathcal{C}_{c}(\mathrm{T};\mathrm{B})\times_{\mathrm{B}}\mathrm{E})\times\mathrm{T}\to\mathrm{B}\) defined by \(\psi((f,x),u)=f(u)\) is continuous (TG, X, p. 28, cor. 1). By Corollary 3 above, the map \(\psi\) has a unique continuous lift \(\theta\colon(\mathcal{C}_{c}(\mathrm{T};\mathrm{B})\times_{\mathrm{B}}\mathrm{E})\times\mathrm{T}\to\mathrm{E}\) such that \(\theta((f,x),t)=x\) for \((f,x)\in\mathcal{C}_{c}(\mathrm{T};\mathrm{B})\times_{\mathrm{B}}\mathrm{E}\) . One thus has \(\theta((f,x),u)=\varphi^{-1}(f,x)(u)\) for \((f,x)\in\mathcal{C}_{c}(\mathrm{T};\mathrm{B})\times_{\mathrm{B}}\mathrm{E}\) and \(u\in\mathrm{T}\) . By Theorem 3 of TG, X, p. 28, the map \((f,x)\mapsto\theta((f,x),\cdot)\) from \(\mathcal{C}_{c}(\mathrm{T};\mathrm{B})\times_{\mathrm{B}}\mathrm{E}\) to \(\mathcal{C}_{c}(\mathrm{T};\mathrm{E})\) is continuous, that is, \(\varphi^{-1}\) is continuous.

Thus, the map \(\varphi\) is a homeomorphism from \(\mathcal{C}_c(\mathrm{T};\mathrm{E})\) onto the fibered product \(\mathcal{C}_c(\mathrm{T};\mathrm{B})\times_{\mathrm{B}}\mathrm{E}\) , whence the proposition (I, p. 8, prop. 2).

